\documentclass[12pt,a4paper]{ctexbook}
\usepackage{geometry}
\geometry{margin=2.2cm}
\usepackage{setspace}
\setstretch{1.35}
\usepackage{hyperref}
\title{花间集·huajianji-x-juan}
\author{古典文献整理}
\date{}
\begin{document}
\maketitle
\tableofcontents
\newpage
\section{河满子·寂寞芳菲暗度}
\textbf{毛熙震}\par\vspace{0.5em}
寂寞芳菲暗度，岁华如箭堪惊。\par
缅想旧欢多少事，转添春思难平。\par
曲槛丝垂金柳，小窗弦断银筝。\par
深院空闻燕语，满园闲落花轻。\par
一片相思休不得，忍教长日愁生。\par
谁见夕陽孤梦，觉来无限伤情。\par

\section{河满子·无语残妆淡薄}
\textbf{毛熙震}\par\vspace{0.5em}
无语残妆淡薄，含羞亸袂轻盈。\par
几度香闺眠过晓，绮窗疏日微明。\par
云母帐中偷惜，水晶枕上初惊。\par
笑靥嫩疑花坼，愁眉翠敛山横。\par
相望只教添怅恨，整鬟时见纤琼。\par
独倚朱扉闲立，谁知别有深情。\par

\section{小重山·梁燕双飞画阁前}
\textbf{毛熙震}\par\vspace{0.5em}
梁燕双飞画阁前，寂寥多少恨，懒孤眠。\par
晓来闲处想君怜，红罗帐，金鸭冷沉烟。\par
谁信损婵娟，倚屏啼玉箸，湿香钿。\par
四肢无力上秋千，群花谢，愁对艳陽天。\par

\section{定西番·苍翠浓陰满院}
\textbf{毛熙震}\par\vspace{0.5em}
苍翠浓陰满院，莺对语，蝶交飞，戏蔷薇。\par
斜日倚栏风好，余香出绣衣。\par
未得玉郎消息，几时归。\par

\section{木兰花·掩朱扉}
\textbf{毛熙震}\par\vspace{0.5em}
掩朱扉，钩翠箔，满院莺声春寂寞。\par
匀粉泪，恨檀郎，一去不归花又落。\par
对斜晖，临小阁，前事岂堪重想着。\par
金带冷，画屏幽，宝帐慵熏兰麝薄。\par

\section{后庭花·莺啼燕语芳菲节}
\textbf{毛熙震}\par\vspace{0.5em}
莺啼燕语芳菲节，瑞庭花发。\par
昔时欢宴歌声揭，管弦清越。\par
自从陵谷追游歇，画梁尘黦。\par
伤心一片如圭月，闲锁宫阙。\par

\section{后庭花·轻盈舞妓含芳艳}
\textbf{毛熙震}\par\vspace{0.5em}
轻盈舞妓含芳艳，竞妆新脸。\par
步摇珠翠修蛾敛，腻鬟云染。\par
歌声慢发开檀点，绣衫斜掩。\par
时将纤手匀红脸，笑招金靥。\par

\section{后庭花·越罗小袖新香蒨}
\textbf{毛熙震}\par\vspace{0.5em}
越罗小袖新香蒨，薄笼金钏。\par
倚栏无语摇轻扇，半遮匀面。\par
春残日暖莺娇懒，满庭花片。\par
怎不教人长相见，画堂深院。\par

\section{酒泉子·闲卧绣纬}
\textbf{毛熙震}\par\vspace{0.5em}
闲卧绣纬，慵想万般情宠。\par
锦檀偏，翘股重，翠云欹。\par
暮天屏上春山碧，映香烟雾隔。\par
蕙兰心，魂梦役，敛蛾眉。\par

\section{酒泉子·钿匣舞鸾}
\textbf{毛熙震}\par\vspace{0.5em}
钿匣舞鸾，隐映艳红修碧。\par
月梳斜，云鬓腻，粉香寒。\par
晓花微敛轻呵展，袅钗金燕软。\par
日初升，帘半卷，对妆残。\par

\section{菩萨蛮·梨花满院飘香雪}
\textbf{毛熙震}\par\vspace{0.5em}
梨花满院飘香雪，高楼夜静风筝咽。\par
斜月照帘帷，忆君和梦稀。\par
小窗灯影背，燕语惊愁态。\par
屏掩断香飞，行云山外归。\par

\section{菩萨蛮·绣帘高轴临塘看}
\textbf{毛熙震}\par\vspace{0.5em}
绣帘高轴临塘看，雨翻荷芰真珠散。\par
残暑晚初凉，轻风渡水香。\par
无聊悲往事，怎奈牵情思。\par
光景暗相催，等闲秋又来。\par

\section{菩萨蛮·天含残碧融春色}
\textbf{毛熙震}\par\vspace{0.5em}
天含残碧融春色，五陵薄幸无消息。\par
尽日掩朱门，离愁暗断魂。\par
莺啼芳树暖，燕拂回塘满。\par
寂寞对屏山，相思醉梦间。\par

\section{浣溪沙·入夏偏宜淡薄妆}
\textbf{李珣}\par\vspace{0.5em}
入夏偏宜淡薄妆，越罗衣褪郁金黄，翠钿檀注助容光。\par
相见无言还有恨，几回拚却又思量，月窗香径梦悠飏。\par

\section{浣溪沙·晚出闲庭看海棠}
\textbf{李珣}\par\vspace{0.5em}
晚出闲庭看海棠，风流学得内家妆，小钗横戴一枝芳。\par
镂玉梳斜云鬓腻，缕金衣透雪肌香，暗思何事立残陽。\par

\section{浣溪沙·访旧伤离欲断魂}
\textbf{李珣}\par\vspace{0.5em}
访旧伤离欲断魂，无因重见玉楼人，六街微雨镂香尘。\par
早为不逢巫峡梦，那堪虚度锦江春，遇花倾酒莫辞频。\par

\section{浣溪沙·红藕花香到槛频}
\textbf{李珣}\par\vspace{0.5em}
红藕花香到槛频，可堪闲忆似花人，旧欢如梦绝音尘。\par
翠叠画屏山隐隐，冷铺纹簟水潾潾，断魂何处一蝉新。\par

\section{渔歌子·楚山青}
\textbf{李珣}\par\vspace{0.5em}
楚山青，湘水绿，春风淡荡看不足。\par
草芊芊，花簇簇，渔艇棹歌相续。\par
信浮沉，无管束，钓回乘月归湾曲。\par
酒盈樽，云满屋，不见人间荣辱。\par

\section{渔歌子·荻花秋}
\textbf{李珣}\par\vspace{0.5em}
荻花秋，潇湘夜，橘洲佳景如屏画。\par
碧烟中，明月下，小艇垂纶初罢。\par
水为乡，篷作舍，鱼羹稻饭常餐也。\par
酒盈杯，书满架，名利不将心挂。\par

\section{渔歌子·柳垂丝}
\textbf{李珣}\par\vspace{0.5em}
柳垂丝，花满树，莺啼楚岸春山暮。\par
棹轻舟，出深浦，缓唱渔歌归去。\par
罢垂纶，还酌醑，孤村遥指云遮处。\par
下长汀，临浅渡，惊起一行沙鹭。\par

\section{渔歌子·九疑山}
\textbf{李珣}\par\vspace{0.5em}
九疑山，三湘水，芦花时节秋风起。\par
水云间，山月里，棹月穿云游戏。\par
鼓清琴，倾绿蚁，扁舟自得逍遥志。\par
任东西，无定止，不议人间醒醉。\par

\section{巫山一段云·有客经巫峡}
\textbf{李珣}\par\vspace{0.5em}
有客经巫峡，停桡向水湄。\par
楚王曾此梦瑶姬，一梦杳无期。\par
尘暗珠帘卷，香消翠帷垂。\par
西风回首不胜悲，暮雨洒空祠。\par

\section{巫山一段云·古庙依青嶂}
\textbf{李珣}\par\vspace{0.5em}
古庙依青嶂，行宫枕碧流。\par
水声山色锁妆楼，往事思悠悠。\par
云雨朝还暮，烟花春复秋。\par
啼猿何必近孤舟，行客自多愁。\par

\section{临江仙·帘卷池心小阁虚}
\textbf{李珣}\par\vspace{0.5em}
帘卷池心小阁虚，暂凉闲步徐徐。\par
芰荷经雨半凋疏，拂堤垂柳，蝉噪夕陽余。\par
不语低鬟幽思远，玉钗斜坠双鱼。\par
几回偷看寄来书，离情别恨，相隔欲何如。\par

\section{临江仙·莺报帘前暖日红}
\textbf{李珣}\par\vspace{0.5em}
莺报帘前暖日红，玉炉残麝犹浓。\par
起来闺思尚疏慵，别愁春梦，谁解此情悰。\par
强整娇姿临宝镜，小池一朵芙蓉。\par
旧欢无处再寻踪，更堪回顾，屏画九疑峰。\par

\section{南乡子·烟漠漠}
\textbf{李珣}\par\vspace{0.5em}
烟漠漠，雨凄凄，岸花零落鹧鸪啼。\par
远客扁舟临野渡，思乡处，潮退水平春色暮。\par

\section{南乡子·兰棹举}
\textbf{李珣}\par\vspace{0.5em}
兰棹举，水纹开，竞携藤笼采莲来。\par
回塘深处遥相见，邀同宴，绿酒一卮红上面。\par

\section{南乡子·归路近}
\textbf{李珣}\par\vspace{0.5em}
归路近，扣舷歌，采真珠处水风多。\par
曲岸小桥山月过，烟深锁，豆蔻花垂千万朵。\par

\section{南乡子·乘彩舫}
\textbf{李珣}\par\vspace{0.5em}
乘彩舫，过莲塘，棹歌惊起睡鸳鸯。\par
游女带香偎伴笑，争窈窕，竞折团荷遮晚照。\par

\section{南乡子·倾绿蚁}
\textbf{李珣}\par\vspace{0.5em}
倾绿蚁，泛红螺，闲邀女伴簇笙歌。\par
避暑信船轻浪里，闲游戏，夹岸荔枝红蘸水。\par

\section{南乡子·云带雨}
\textbf{李珣}\par\vspace{0.5em}
云带雨，浪迎风，钓翁回棹碧湾中。\par
春酒香熟鲈鱼美，谁同醉？\par
缆却扁舟篷底睡。\par

\section{南乡子·沙月静}
\textbf{李珣}\par\vspace{0.5em}
沙月静，水烟轻，芰荷香里夜船行。\par
绿鬟红脸谁家女？\par
遥相顾，缓唱棹歌极浦去。\par

\section{南乡子·渔市散}
\textbf{李珣}\par\vspace{0.5em}
渔市散，渡船稀，越南云树望中微。\par
行客待湘天欲暮，送春浦，愁听猩猩啼瘴雨。\par

\section{南乡子·拢云髻}
\textbf{李珣}\par\vspace{0.5em}
拢云髻，背犀梳，焦红衫映绿罗裙。\par
越王台下春风暖，药盈岸，游赏每邀邻女伴。\par

\section{南乡子·相见处}
\textbf{李珣}\par\vspace{0.5em}
相见处，晚晴天，斜桐花下越台前。\par
暗里回眸深属意，遗双翠，骑象背人先过水。\par

\section{女冠子·星高月午}
\textbf{李珣}\par\vspace{0.5em}
星高月午，丹桂青松深处。\par
醮坛开，金磬敲清露，珠幢立翠苔。\par
步虚声缥缈，想象思徘徊。\par
晓天归去路，指蓬莱。\par

\section{女冠子·春山夜静}
\textbf{李珣}\par\vspace{0.5em}
春山夜静，愁闻洞天疏磬。\par
玉堂虚，细雾垂珠佩，轻烟曳翠裙。\par
对花情脉脉，望月步徐徐。\par
刘阮今何处？绝来书。\par

\section{酒泉子·寂寞青楼}
\textbf{李珣}\par\vspace{0.5em}
寂寞青楼，风触绣帘珠碎撼。\par
月朦胧，花暗淡，锁春愁。\par
寻想往事依稀梦，泪脸桃红色重。\par
鬓欹蝉。钗坠凤，思悠悠。\par

\section{酒泉子·雨渍花零}
\textbf{李珣}\par\vspace{0.5em}
雨渍花零，红散香凋池两岸。\par
别情遥，春歌断，掩银屏。\par
孤帆早晚离三楚，闲理钿筝愁几许？\par
曲中情，弦上语，不堪听。\par

\section{酒泉子·秋雨联绵}
\textbf{李珣}\par\vspace{0.5em}
秋雨联绵，声散败荷丛里。\par
那堪深夜枕前听，酒初醒。\par
牵愁惹思更无停，烛暗香凝天欲曙。\par
细和烟，冷和雨，透帘旌。\par

\section{酒泉子·秋月婵娟}
\textbf{李珣}\par\vspace{0.5em}
秋月婵娟，皎洁碧纱窗外。\par
照花穿竹冷沉沉，印池心。\par
凝露滴，砌蛩吟，惊觉谢娘残梦。\par
夜深斜傍枕前来，影徘徊。\par

\section{望远行·春日迟迟思寂寥}
\textbf{李珣}\par\vspace{0.5em}
春日迟迟思寂寥，行客关山路遥。\par
琼窗时听语莺娇，柳丝牵恨一条条。\par
休晕绣，罢吹萧，貌逐残花暗凋。\par
同心犹结旧裙腰，忍辜风月度良宵。\par

\section{望远行·露滴幽庭落叶时}
\textbf{李珣}\par\vspace{0.5em}
露滴幽庭落叶时，愁聚萧娘柳眉。\par
玉郎一去负佳期，水云迢递雁书迟。\par
屏半掩，枕斜欹，蜡泪无言对垂。\par
吟蛩断续漏频移，入窗明月鉴空帷。\par

\section{菩萨蛮·回塘起波纹细}
\textbf{李珣}\par\vspace{0.5em}
回塘起波纹细，剌桐花里门斜闭。\par
残日照平芜，双双飞鹧鸪。\par
征帆何处客，相见还相隔。\par
不语欲魂消，望中烟水遥。\par

\section{菩萨蛮·等闲将度三春景}
\textbf{李珣}\par\vspace{0.5em}
等闲将度三春景，帘垂碧砌参差影。\par
曲槛日初斜，杜鹃啼落花。\par
恨君容易处，又话潇湘去。\par
凝思倚屏山，泪流红脸斑。\par

\section{菩萨蛮·隔帘微雨双飞燕}
\textbf{李珣}\par\vspace{0.5em}
隔帘微雨双飞燕，砌花零落红深浅。\par
捻得宝筝凋，心随征棹遥。\par
楚天云外路，动便经年去。\par
香断画屏深，旧欢何处寻。\par

\section{西溪子·金缕翠钿浮动}
\textbf{李珣}\par\vspace{0.5em}
金缕翠钿浮动，妆罢小窗圆梦。\par
日高时，春已老，人来到，满地落花慵扫。\par
无语倚屏风，泣残红。\par

\section{虞美人·金笼莺报天将曙}
\textbf{李珣}\par\vspace{0.5em}
金笼莺报天将曙，惊起分飞处。\par
夜来潜与玉郎期，多情不觉酒醒迟，失归期。\par
映花避月遥相送，腻髻偏凤。\par
欲回娇步入香闺，倚屏无语捻云篦，翠眉低。\par

\section{河传·去去}
\textbf{李珣}\par\vspace{0.5em}
去去，何处？迢迢巴楚，山水相连。\par
朝云暮雨，依旧十二峰前，猿声到客船。\par
愁肠岂异丁香结？因离别，故国音书绝。\par
想佳人花下，对明月春风，恨应同。\par

\section{河传·春暮}
\textbf{李珣}\par\vspace{0.5em}
春暮，微雨。送君南浦，愁敛双蛾。\par
落花深处，啼鸟似逐离散，粉檀珠泪和。\par
临流更把同心结，情哽咽，后会何时节？\par
不堪回首，相望已隔汀洲，檐声幽。\par

\end{document}